% !TEX root = ../main.tex
%% Sujet type bac n° 3 — sujet inédit OptiTAB (valeurs vérifiées : verifications/sujets_bac.py)

\sujetBac{3}{{L'arôme de banane : synthèse, équilibre, spectroscopie}/8,{Un satellite d'observation : gravitation, diffraction}/7,{Du thé à la bonne température : énergie, transferts thermiques}/5}

\begin{exoBac}{8}{L'arôme de banane}
L'éthanoate de 3-méthylbutyle, appelé acétate d'isoamyle, donne à la banane son odeur caractéristique. Les industriels
de l'agroalimentaire le synthétisent par estérification, selon une transformation modélisée par la réaction
\[\ce{CH3COOH + C5H11OH <=> CH3COOC5H11 + H2O}\qquad K=4{,}0.\]
\begin{center}
\setchemfig{atom sep=1.9em}
\chemfig{-[:30]-[:-30](-[:-90])-[:30]-[:-30]-[:30]OH}\qquad\qquad
\chemfig{-[:30](=[:90]O)-[:-30]O-[:30]-[:-30]-[:30](-[:90])-[:-30]}\par\smallskip
{\legende{3-méthylbutan-1-ol (à gauche) et éthanoate de 3-méthylbutyle (à droite).}}
\end{center}
\begin{docu}{1}{Le protocole}
Dans un ballon, on introduit \qty{0,200}{mol} de 3-méthylbutan-1-ol, \qty{0,400}{mol} d'acide éthanoïque, quelques gouttes d'acide
sulfurique concentré et de la pierre ponce. On chauffe à reflux pendant une heure. Après refroidissement, on lave la phase
organique avec une solution d'hydrogénocarbonate de sodium, $\ce{Na+}+\ce{HCO3-}$ : une effervescence se produit. On sèche,
puis on distille : on recueille \qty{20,8}{g} d'ester.
\end{docu}
\begin{docu}{2}{Données physico-chimiques}
{\small\begin{tabular}{@{}lccc@{}}
Espèce & 3-méthylbutan-1-ol & acide éthanoïque & ester \\
Masse molaire (\unit{g.mol^{-1}}) & 88,1 & 60,1 & 130,2 \\
Masse volumique (\unit{g.mL^{-1}}) & 0,81 & 1,05 & 0,87 \\
Température d'ébullition (\unit{\degreeCelsius}) & 131 & 118 & 142 \\
\end{tabular}}\par\smallskip
Électronégativités : H 2,2 ; C 2,6 ; O 3,4. Couple $\ce{CO2},\ce{H2O}/\ce{HCO3-}$. En infrarouge, la liaison O--H d'un alcool absorbe
vers 3200 -- \qty{3400}{cm^{-1}} (bande large) et la liaison C$=$O d'un ester vers 1735 -- \qty{1750}{cm^{-1}}.
\end{docu}
\schema{bac-s3-ir}
\begin{questions}
\item Recopier les deux formules topologiques, entourer les groupes caractéristiques et nommer les familles fonctionnelles.
\item Calculer le volume d'acide éthanoïque à prélever.
\item Quel est l'intérêt du chauffage à reflux ? Quel est le rôle de l'acide sulfurique ?
\item Calculer le rendement de la synthèse.
\item Exprimer la constante d'équilibre en fonction de l'avancement final $x_\mathrm{f}$, puis vérifier que $x_\mathrm{f}=\qty{0,169}{mol}$. En
déduire le rendement maximal que permet l'équilibre, et proposer une explication à l'écart avec le rendement obtenu.
\item \initiative Montrer qu'avec un mélange initial de \qty{0,200}{mol} de chaque réactif, le rendement à l'équilibre ne serait
que de 67~\%. Pourquoi l'acide éthanoïque est-il choisi en excès plutôt que l'alcool ?
\item Écrire l'équation de la réaction responsable de l'effervescence observée lors du lavage. Quel est l'intérêt de ce
lavage ?
\item Attribuer chaque spectre infrarouge à l'espèce correspondante, en justifiant.
\item Dans une étape du mécanisme, l'atome d'oxygène du groupe hydroxyle de l'alcool forme une liaison avec l'atome de
carbone du groupe carboxyle, préalablement protoné, de l'acide. Identifier le site donneur et le site accepteur de doublet
d'électrons, puis justifier le rôle de chacun.
\end{questions}
\end{exoBac}

\begin{exoBac}{7}{Un satellite d'observation}
Les satellites Pléiades photographient la surface terrestre avec une résolution inférieure au mètre. Ils décrivent, dans
le référentiel géocentrique supposé galiléen, une orbite circulaire à l'altitude $h=\qty{694}{km}$.
\begin{donnees}
$G=\qty{6,67e-11}{N.m^2.kg^{-2}}$ ; masse de la Terre $M_\mathrm{T}=\qty{5,97e24}{kg}$ ; rayon de la Terre $R_\mathrm{T}=\qty{6371}{km}$ ; période de rotation
de la Terre dans le référentiel géocentrique : \qty{86164}{s}. Pour un instrument d'ouverture $D$, la diffraction limite la
plus petite séparation angulaire observable à $\theta\approx\dfrac{\lambda}{D}$.
\end{donnees}
\begin{questions}
\item Définir le référentiel géocentrique. Donner l'expression vectorielle de la force exercée par la Terre sur le satellite,
de masse $m$, en utilisant le vecteur unitaire $\vec u_n$ de la base de Frenet.
\item Montrer, à l'aide de la deuxième loi de Newton, que le mouvement est uniforme et que la vitesse du satellite vaut
$v=\sqrt{\dfrac{G\,M_\mathrm{T}}{R_\mathrm{T}+h}}$.
\item Calculer cette vitesse.
\item Établir l'expression de la période de révolution, la calculer, puis déterminer le nombre de tours effectués par jour.
\item Montrer que le rapport $\dfrac{T^2}{r^3}$, où $r$ est le rayon de l'orbite, ne dépend que de la masse de la Terre. Quel nom donne-t-on
à ce résultat ?
\item Calculer l'accélération du satellite. Un objet à bord paraît flotter : est-ce parce que la pesanteur y est nulle ?
\item Le télescope du satellite a un miroir de diamètre $D=\qty{0,65}{m}$ et observe à $\lambda=\qty{600}{nm}$. Calculer la plus petite
séparation angulaire observable, puis la taille du plus petit détail visible au sol.
\item \initiative Un satellite géostationnaire observerait la même région en permanence. Avec le même télescope, quelle
serait la taille du plus petit détail visible ? Conclure.
\end{questions}
\end{exoBac}

\begin{exoBac}{5}{Du thé à la bonne température}
Pour préparer un thé, on chauffe \qty{1,0}{L} d'eau dans une bouilloire électrique de puissance \qty{2200}{W} : l'eau passe de \qty{18}{\degreeCelsius} à
\qty{100}{\degreeCelsius} en \qty{3,0}{min}. Le thé infuse, puis refroidit dans une tasse posée dans une pièce à \qty{20}{\degreeCelsius}.
\begin{donnees}
Capacité thermique massique de l'eau : $c=\qty{4180}{J.kg^{-1}.K^{-1}}$ ; masse volumique de l'eau : \qty{1,0}{kg.L^{-1}}. Loi de Newton :
la température $T$ du thé vérifie $\dfrac{\mathrm{d}T}{\mathrm{d}t}=-\dfrac{T-T_\mathrm{ext}}{\tau}$, où $\tau$ est une constante.
\end{donnees}
\schema{bac-s3-the}
\begin{questions}
\item Calculer la variation d'énergie interne de l'eau pendant le chauffage. Quel mode de transfert d'énergie l'a
provoquée ?
\item Calculer l'énergie électrique consommée par la bouilloire, puis son rendement. Proposer deux causes de pertes.
\item Vérifier que $T(t)=T_\mathrm{ext}+\left(T_0-T_\mathrm{ext}\right)\mathrm{e}^{-t/\tau}$ est solution de l'équation de Newton.
\item À l'aide de la figure, montrer que $\tau\approx\qty{25}{min}$.
\item \initiative Au bout de combien de temps le thé atteint-il \qty{40}{\degreeCelsius} ? Comment modifier la tasse pour que le thé refroidisse
moins vite ?
\end{questions}
\end{exoBac}
